\section{Related work}\label{sec:related}

\paragraph{Saturation of static LLM benchmarks.}
By May 2026, the headline static-resettable benchmarks are at or past
saturation. MMLU and MMLU-Pro sit above 88\,\%; GPQA-Diamond above
94\,\%; AIME 2025/26 and MATH-500 above 91\,\%; SWE-bench Verified at
93.9\,\% (Claude Mythos)~\cite{swebench}; OSWorld-Verified has been
crossed by Simular Agent S2 and GPT-5.4~\cite{osworld}. HLE rose from
8\,\% (Jan 2025, o1) to 65\,\% (Claude Mythos, May 2026)~\cite{hle};
ARC-AGI-2 from 53\,\% (Dec 2025) to 85\,\% (GPT-5.5, May 2026)
~\cite{arcagi2}. The capabilities that did \emph{not} saturate are all
about calibrated, persistent agency: ARC-AGI-3 $<$1\,\%~\cite{arcagi3},
SWE-EVO 25\,\%, MemoryAgentBench selective-forgetting
$<$30\,\%~\cite{memoryagentbench}.

\paragraph{Verbal calibration of LLMs.} Kadavath
et~al.~\cite{kadavath2022} showed LLMs have non-trivial verbal
calibration on factual QA; Tian et~al.\ (2023) benchmarked verbalised
probabilities. Recent reasoning-faithfulness work
\cite{anthropic_faithfulness} found that models verbalise hints in only
25--39\,\% of cases and that \emph{unfaithful} chains-of-thought are
systematically \emph{longer} than faithful ones --- a finding directly
relevant to PROPHET's calibration emphasis.

\paragraph{Reliability beyond pass@1.} \cite{beyond_pass1} formalised
the Reliability Decay Curve, Variance Amplification Factor, and Model
Overreach Point. \cite{long_horizon_mirage} found recovery is the
dominant production failure mode. PROPHET operationalises these as
publishable scores.

\paragraph{Proper scoring rules in machine-learning evaluation.}
Brier (1950) introduced the canonical quadratic scoring rule for
probability forecasts; Gneiting and Raftery~\cite{gneiting2007} survey
the theory and practical alternatives. Guo et~al.~\cite{guo2017}
established ECE and temperature scaling as the modern calibration
benchmarks for deep classifiers. We adopt Brier as the calibration term
because it is bounded, smooth, easy to communicate, and admits closed-
form expected payoff at known $P[y=1\!\mid\!t]$.

\paragraph{Agent benchmarks and gaming.} GAIA~\cite{gaia} established
multi-step tool-use as a measurable surface; AgentBench, MMAU, OSWorld
~\cite{osworld} followed. The Berkeley RDI April~2026
audit~\cite{berkeley_rdi_audit} showed trivial exploits gave 100\,\% on
eight major agent benchmarks. We adopt their attack taxonomy as a pre-
publication red-team checklist (Section~\ref{sec:gaming}).

\paragraph{Multiple-comparison correction in benchmark papers.}
Demsar (2006) argued for non-parametric paired tests when comparing
multiple classifiers; Dem\v{s}ar's argument applies a fortiori at our
scale (twelve task families, $\sim$15 agents, three primary metrics).
We default to Holm--Bonferroni for family-wise error rate; results are
also reported under Benjamini--Hochberg FDR.

\paragraph{Concurrent benchmarks.} \cite{memoryagentbench} measures
selective forgetting; \cite{beyond_pass1} measures reliability decay;
PROPHET is complementary: it measures calibrated commitment under a
proper scoring rule across a heterogeneous task suite. None of the
existing benchmarks publishes a Pareto frontier on
$(\text{net payoff}, \mathrm{ECE})$.
